% !TEX root = ../../cookbook.tex
\newpage
\recipe{Soupe à l'Oignon}{Pour 4 portions}{}{}

\begin{multicols}{2}
	\subsection*{Ingrédients}
	\subsubsection*{Soupe}
	\begin{ingredients}
		\item 6 oignons jaunes
		\item 55 g de beurre
		\item 125 ml de vin rouge
		\item 30 ml de cognac
		\item \sfrac{1}{2} L de bouillon de poulet
		\item \sfrac{1}{2} L de bouillon de boeuf
		\item 1 c.à.s de farine
		\item 1 pincée de muscade
	\end{ingredients}
	
	\columnbreak
	
	\subsection*{}
	\subsubsection*{Garniture}
	\begin{ingredients}
		\item 200 g de Gruyére rapé
		\item Tranches de baguette grillées
		\item 1 gousse d'ail
	\end{ingredients}
	\vfill
\end{multicols}

\medskip

% --- Preparation ---
Couper les \textbf{oignons} en quartiers, puis les émincer pour obtenir des courtes lamières. 
Dans une grande casserole, à feu moyen, attendrir les oignons dans le \textbf{beurre} pendant 15 minutes. 
Poursuivre à few élevé 15 minutes, ou jusqu'à ce que les oignons soient caramélisées en remuant constamment et en grattant le fond. 
Déglacer avec le \textbf{vin} et le \textbf{cognac}, et réduire presque sec. 
Ajouter les \textbf{bouillons}, la \textbf{farine}, et la \textbf{muscade}.
Mijoter à feu moien-élevé 30 minutes, ou jusqu'à ce que la soupe ait réduit de moitié.
Saler et poivrer. \\

Placer les bols de soupe dans le four avec le pain frotté avec l'ail, bien couvrir avec le fromage et griller. 


\vspace*{0.75 cm}

% --- Description ---
\begin{recipeblock}
	
\end{recipeblock}